% =====================================================================
% CHAPTER 4: PROJECT DEVELOPMENT
% =====================================================================
\chapter{\dev}
\label{ch:development}

% Write the introduction to the development here...

% ---------------------------------------------------------
\section{\sectReq}
\label{sec:requirements}
{\small\itshape\color{gray} \reqHelpText}

% Write the requirements analysis here...

\subsection{\sectFunc}
\label{subsec:functional_requirements}
% Example table of functional requirements (replace it with your own)
\exFRIntro

\begin{table}[H]
    \centering
    \small
    \begin{tabular}{|p{1.3cm}|p{3.6cm}|p{5.9cm}|p{2.4cm}|}
        \hline
        \textbf{\exHeadID} & \textbf{\exHeadName} & \textbf{\exHeadDescription} & \textbf{\exHeadPriority} \\ \hline
        \exFRPrefix-01 & \exFRUsers  & \exFRUsersDesc  & \exHigh \\ \hline
        \exFRPrefix-02 & \exFRIngest & \exFRIngestDesc & \exCritical \\ \hline
        \exFRPrefix-03 & \exFRAlerts & \exFRAlertsDesc & \exMedium \\ \hline
    \end{tabular}
    \caption{\exFRCaption}
    \label{tab:functional_requirements}
\end{table}

\subsection{\sectCases}
\label{subsec:use_cases}
% Example use case diagram drawn with TikZ (replace it with your own)
\exUseCaseIntro

\begin{figure}[H]
    \centering
    \begin{tikzpicture}[
        usecase/.style={ellipse, draw, thick, fill=blue!5, minimum width=3.0cm, minimum height=0.9cm, text width=2.3cm, align=center, font=\small, inner sep=1pt},
        assoc/.style={thin},
        rel/.style={-{Stealth[length=2mm]}, dashed},
        stereo/.style={font=\scriptsize\itshape, fill=white, inner sep=1pt},
        actor/.pic={
            \draw[thick] (0,0.55) circle (0.18);
            \draw[thick] (0,0.37) -- (0,-0.15);
            \draw[thick] (-0.3,0.2) -- (0.3,0.2);
            \draw[thick] (0,-0.15) -- (-0.22,-0.55);
            \draw[thick] (0,-0.15) -- (0.22,-0.55);
        }
    ]
    % System boundary
    \draw[thick, fill=gray!5, rounded corners=4pt] (-4.4,-3.2) rectangle (4.4,3.6);
    \node[font=\bfseries] at (0,3.2) {\exUCSystem};

    % Use cases
    \node[usecase] (login)   at (-2.3, 2.2)  {\exUCLogin};
    \node[usecase] (cred)    at ( 2.3, 2.2)  {\exUCCredentials};
    \node[usecase] (create)  at (-2.3, 0.7)  {\exUCCreate};
    \node[usecase] (attach)  at ( 2.3, 0.7)  {\exUCAttach};
    \node[usecase] (status)  at (-2.3,-0.8)  {\exUCStatus};
    \node[usecase] (resolve) at ( 2.3,-0.8)  {\exUCResolve};
    \node[usecase] (users)   at ( 0.0,-2.3)  {\exUCUsers};

    % Actors
    \pic at (-6.2, 0.7) {actor};
    \node[font=\small] at (-6.2,-0.05) {\exUCCustomer};
    \pic at ( 6.2,-0.3) {actor};
    \node[font=\small] at ( 6.2,-1.05) {\exUCTechnician};
    \pic at ( 6.2,-2.5) {actor};
    \node[font=\small] at ( 6.2,-3.25) {\exUCAdmin};

    % Associations
    \draw[assoc] (-5.8, 0.8) -- (login);
    \draw[assoc] (-5.8, 0.8) -- (create);
    \draw[assoc] (-5.8, 0.8) -- (status);
    \draw[assoc] ( 5.8,-0.2) -- (resolve);
    \draw[assoc] ( 5.8,-2.4) -- (users);

    % Include and extend relationships
    \draw[rel] (login) -- node[stereo, above=1pt] {$\ll$include$\gg$} (cred);
    \draw[rel] (attach) -- node[stereo, above=1pt] {$\ll$extend$\gg$} (create);
    \end{tikzpicture}
    \caption{\exUseCaseCaption}
    \label{fig:use_case_diagram}
\end{figure}

\subsection{\sectNonFunc}
\label{subsec:nonfunctional_requirements}
% Example table of non-functional requirements (replace it with your own)
\exNFRIntro

\begin{table}[H]
    \centering
    \small
    \begin{tabular}{|p{1.5cm}|p{3.1cm}|p{6.1cm}|p{2.4cm}|}
        \hline
        \textbf{\exHeadID} & \textbf{\exHeadAttribute} & \textbf{\exHeadDescription} & \textbf{\exHeadPriority} \\ \hline
        \exNFRPrefix-01 & \exNFRSecurity     & \exNFRSecurityDesc     & \exHigh \\ \hline
        \exNFRPrefix-02 & \exNFRPerformance  & \exNFRPerformanceDesc  & \exMedium \\ \hline
        \exNFRPrefix-03 & \exNFRAvailability & \exNFRAvailabilityDesc & \exHigh \\ \hline
        \exNFRPrefix-04 & \exNFRScalability  & \exNFRScalabilityDesc  & \exMedium \\ \hline
    \end{tabular}
    \caption{\exNFRCaption}
    \label{tab:nonfunctional_requirements}
\end{table}

% ---------------------------------------------------------
\section{\sectArchDesign}
\label{sec:design}
{\small\itshape\color{gray} \archDesignHelpText}

% Write the design and architecture here...

\subsection{\sectArch}
\label{subsec:architecture}
% Example architecture diagram drawn with TikZ (replace it with your own)
\begin{figure}[H]
    \centering
    \begin{tikzpicture}[
        node distance=1.5cm,
        block/.style={rectangle, draw, fill=blue!5, text width=6cm, text centered, rounded corners, minimum height=1.2cm, font=\bfseries},
        arrow/.style={thick,<->,>=stealth, font=\scriptsize\itshape}
    ]
    % Layers
    \node (pres) [block] {\exArchPresentation \\ \textnormal{\small (Frontend / UI)}};
    \node (logic) [block, below=of pres] {\exArchBusiness \\ \textnormal{\small (\exArchBusinessSub)}};
    \node (data) [block, below=of logic] {\exArchData \\ \textnormal{\small (\exArchDataSub)}};

    % Communication between layers
    \draw [arrow] (pres.south) -- node[right] {\exArchRequests} (logic.north);
    \draw [arrow] (logic.south) -- node[right] {\exArchQueries} (data.north);
    \end{tikzpicture}
    \caption{\exArchCaption}
    \label{fig:architecture}
\end{figure}

\subsection{\sectDB}
\label{subsec:database}
% Write the database design here...

\subsubsection{\sectUML}
% Example entity-relationship model drawn with TikZ (replace it with your own)
\begin{figure}[H]
    \centering
    \begin{tikzpicture}[
        node distance=2cm,
        entity/.style={rectangle, draw, thick, fill=gray!5, text width=3.2cm, text centered, minimum height=0.8cm, font=\bfseries\small},
        attr/.style={ellipse, draw, fill=white, text width=1.6cm, text centered, font=\tiny, inner sep=1pt},
        rel/.style={diamond, draw, thick, fill=white, text width=1.2cm, text centered, font=\scriptsize\itshape, inner sep=0pt}
    ]
    % 1. Entities
    \node (user) [entity] {\exERUser};
    \node (ticket) [at={(7,0)}, entity] {\exERTicket};

    % 2. Relationship
    \node (rel) [rel] at (3.5,0) {\exEROwns};

    % 3. Attributes
    \node (u_id) [attr] at (-1.5, 1.2) {\underline{id\_user}};
    \node (u_email) [attr] at (-1.5, -1.2) {email};
    \node (t_id) [attr] at (8.5, 1.2) {\underline{id\_ticket}};
    \node (t_subject) [attr] at (8.5, -1.2) {\exERSubject};

    % 4. Connections and cardinality
    \draw [thick] (user) -- (rel) node[midway, above, font=\tiny] {1};
    \draw [thick] (rel) -- (ticket) node[midway, above, font=\tiny] {N};

    % Attribute lines
    \draw [thin] (user) -- (u_id);
    \draw [thin] (user) -- (u_email);
    \draw [thin] (ticket) -- (t_id);
    \draw [thin] (ticket) -- (t_subject);
    \end{tikzpicture}
    \caption{\exERCaption}
    \label{fig:er_model}
\end{figure}

\subsubsection{\exDictTitle}
% Example data dictionary (replace it with your own)
\exDictIntro

\begin{table}[H]
    \centering
    \small
    \begin{tabular}{|l|l|l|p{6cm}|}
        \hline
        \textbf{\exHeadField} & \textbf{\exHeadType} & \textbf{\exHeadKey} & \textbf{\exHeadDescription} \\ \hline
        id\_user & INT          & PK & \exDictId \\ \hline
        email    & VARCHAR(100) & UK & \exDictEmail \\ \hline
        password & VARCHAR(255) & -- & \exDictPassword \\ \hline
        role     & ENUM         & -- & \exDictRole \\ \hline
    \end{tabular}
    \caption{\exDictCaption}
    \label{tab:data_dictionary_users}
\end{table}

\subsection{\sectUI}
\label{subsec:user_interface}
% Example wireframe drawn with TikZ (replace it with your own screens)
\begin{figure}[H]
    \centering
    \begin{tikzpicture}[
        window/.style={draw=gray!80, thick, fill=white, minimum width=11.5cm, minimum height=7.5cm, rounded corners=6pt, drop shadow={opacity=0.3}},
        header/.style={draw=none, fill=gray!20, minimum width=11.5cm, minimum height=1cm, rounded corners=6pt},
        addressbar/.style={draw=gray!30, fill=white, minimum width=7cm, minimum height=0.5cm, rounded corners=10pt},
        sidebar/.style={draw=none, fill=gray!10, minimum width=2.5cm, minimum height=6.5cm},
        btn_active/.style={draw=none, fill=orange!80, text=white, font=\tiny\bfseries, rounded corners=3pt, minimum width=2cm, minimum height=0.6cm},
        btn_normal/.style={draw=none, text=gray!80, font=\tiny, minimum width=2cm, minimum height=0.6cm},
        close_dot/.style={circle, minimum size=6pt, inner sep=0pt, fill=orange!80}
    ]
    % 1. Main window with shadow
    \node (browser) [window] {};

    % 2. Top toolbar
    \node (head) [header] at (0, 3.25) {};

    % Window control buttons
    \node [close_dot] at (5.2, 3.4) {};
    \draw[gray!60, thick] (4.8, 3.4) -- (5.0, 3.4); % Minimise
    \draw[gray!60, thick] (4.4, 3.3) rectangle (4.6, 3.5); % Maximise

    % Address bar
    \node (url) [addressbar] at (0, 3.25) {};
    \node at (0, 3.25) {\tiny \color{gray} app.ticketing-system.local};

    % 3. Left sidebar
    \node (side) [sidebar] at (-4.5, -0.5) {};

    % Sidebar buttons
    \node (b1) [btn_active] at (-4.5, 2.0) {\exUIHome};
    \node (b2) [btn_normal] at (-4.5, 1.3) {\exUITickets};
    \node (b3) [btn_normal] at (-4.5, 0.6) {\exUIUsers};
    \node (b4) [btn_normal] at (-4.5, -0.1) {\exUISettings};

    % 4. Main content
    \node[anchor=west] at (-3.0, 2.3) {\large \textbf{\exUITitle}};
    \draw[orange, thick] (-3.0, 2.0) -- (-1.0, 2.0); % Underline

    % Summary cards
    \draw[draw=gray!20, fill=orange!5, rounded corners=4pt] (-3.0, 0.5) rectangle (-0.5, 1.5);
    \node at (-1.75, 1.2) {\tiny \textbf{\exUIPending}};
    \node at (-1.75, 0.8) {\large \color{orange!80!black} 5};

    \draw[draw=gray!20, fill=gray!5, rounded corners=4pt] (0, 0.5) rectangle (2.5, 1.5);
    \node at (1.25, 1.2) {\tiny \textbf{\exUIResolved}};
    \node at (1.25, 0.8) {\large 142};

    % Data table
    \draw[draw=gray!20, fill=white] (-3.0, -3.0) rectangle (5.0, 0);
    \fill[gray!10] (-3.0, -0.4) rectangle (5.0, 0); % Table header
    \node[anchor=west] at (-2.8, -0.2) {\tiny \textbf{\exUIReference}};
    \node[anchor=west] at (-0.5, -0.2) {\tiny \textbf{\exUIUser}};
    \node[anchor=west] at (3.0, -0.2) {\tiny \textbf{\exUIPriority}};

    \node at (1.0, -1.7) {\tiny \color{gray} \exUISync};
    \end{tikzpicture}
    \caption{\exUICaption}
    \label{fig:wireframe}
\end{figure}

% ---------------------------------------------------------
\section{\sectImpl}
\label{sec:implementation}
{\small\itshape\color{gray} \implHelpText}

% Write the implementation here...

% ------------------- Example: algorithm -------------------
\begin{algorithm}[H] % [H] keeps the algorithm exactly here
\caption{\exAlgorithmCaption}\label{alg:example}
\begin{algorithmic}
\Require $n \geq 0$
\Ensure $y = x^n$
\State $y \gets 1$
\State $X \gets x$
\State $N \gets n$
\While{$N \neq 0$}
    \If{$N$ is even}
        \State $X \gets X \times X$
        \State $N \gets \frac{N}{2}$
    \ElsIf{$N$ is odd}
        \State $y \gets y \times X$
        \State $N \gets N - 1$
    \EndIf
\EndWhile
\end{algorithmic}
\end{algorithm}

% ------------------- Example: Java code -------------------
\begin{lstlisting}[language=Java]
public static void main(String args[]) {
    for(int i = 0; i <= 12; i++) { // comment
        System.out.print("12 * "+ i + " = " + 12 * i + "\n");
    }
}
\end{lstlisting}

% ------------------- Example: SQL code --------------------
\begin{lstlisting}[language=SQL,numbers=none]
SELECT *
FROM data
WHERE day BETWEEN 8 AND 10;
\end{lstlisting}
